% Part: sets-functions-relations
% Chapter: arithmetization
% Section: checking-details
%
\documentclass[../../../include/open-logic-section]{subfiles}


\begin{document}
	
\olfileid{sfr}{arith}{check}
\olsection{Geordende ringen en lichamen}
In dit hoofdstuk zeiden we een paar keer dat bepaalde definities
``werken zoals het hoort''. In deze technische bijlage leggen we uit
wat dat betekent. We laten ook in grote lijnen zien hoe je kunt
bewijzen dat de definities echt ``goed'' werken.
In \olref[int]{sec} hebben we vastgelegd hoe je de getallen in $\Int$
optelt en vermenigvuldigt. We moeten nu laten zien dat $\Int$ daardoor
de structuur krijgt die we willen. Om precies te zijn: de gehele
getallen moeten samen een commutatieve ring vormen.
\begin{defn}
	Een \emph{commutatieve ring} is een verzameling $S$ met twee speciale elementen, $0$ en $1$, en twee rekenbewerkingen, $+$ en $\times$. De volgende acht formules moeten dan kloppen:
	\begin{align*}
		\emph{Associativiteit (anders groeperen)}&&a + (b+ c) & = (a + b) + c \\
		&& (a \times b) \times c & = a \times (b\times c)\\
		\emph{Commutativiteit (volgorde omwisselen)}&&a + b &= b+ a  \\
		&&  a \times b&= b\times a\\ 
		\emph{Neutrale elementen (0 bij optellen en 1 bij vermenigvuldigen veranderen niets)}&&a + 0 &= a \\
		&& a \times 1 &= a\\
		\emph{Additieve inverse (tegengestelde bij optellen)}&&(\exists b\in S)0&=a + b\\
		\emph{Distributiviteit (vermenigvuldigen over een som)}&&a \times (b+ c ) &= (a \times b) + (a \times c)
	\end{align*}
	Deze formules moeten kloppen voor elke keuze van de letters uit $S$; dat schrijven we er niet elke keer apart bij. Let ook op: de elementen die hier $0$ en~$1$ heten, hoeven niet de natuurlijke getallen nul en één te zijn.
\end{defn}
Om te bewijzen dat de gehele getallen een commutatieve ring vormen,
moeten we dus deze acht regels nalopen. Elke regel op zichzelf is
eenvoudig, maar samen is het nogal wat werk. Hieronder bewijzen we
bijvoorbeeld de \emph{associativiteit} van het optellen: de uitkomst
verandert niet als je de haakjes anders zet.
\begin{proof} 
Kies $i, j, k \in \Int$. We kunnen dan natuurlijke getallen $a_1, b_1,
a_2, b_2, a_3, b_3 \in \Nat$ vinden zodat $i = \equivrep{a_1, b_1}{}$ en $j =
\equivrep{a_2,b_2}{}$ en $k = \equivrep{a_3, b_3}{}$. (Om dit
makkelijker te lezen, schrijven we ``$\equivrep{x, y}{}$'' in plaats van
``$\equivrep{x, y}{\Intequiv}$''. Dat doen we in dit hele
deel.) Nu rekenen we:
\begin{align*}
	i + (j + k) &= \equivrep{a_1, b_1}{}+(\equivrep{a_2, b_2}{} + \equivrep{a_3, b_3}{}) \\
	&= \equivrep{a_1,  b_1}{} + \equivrep{a_2+a_3, b_2+b_3}{}\\
	&= \equivrep{a_1 + (a_2 + a_3), b_1 + (b_2 + b_3)}{}\\
	&= \equivrep{(a_1 + a_2) + a_3, (b_1 + b_2) + b_3}{}\\
	&= \equivrep{a_1 + a_2, b_1 + b_2}{} + \equivrep{a_3, b_3}{}\\
	&= (\equivrep{a_1, b_1}{} + \equivrep{a_2, b_2}{}) + \equivrep{a_3, b_3}{}\\
	&= (i+j) + k
\end{align*}
Hier gebruiken we gewoon de bekende regels voor optellen op $\Nat$.
\end{proof}
Nu laten we zien dat elk geheel getal bij optellen een tegengestelde
heeft, oftewel een \emph{additieve inverse}:
\begin{proof}
Kies $i \in \Int$ en schrijf $i = \equivrep{a,b}{}$ met $a,b \in
\Nat$. Neem $j = \equivrep{b,a}{} \in
\Int$. Voor natuurlijke getallen geldt $(a+b) + 0 = 0 + (b+a)$.
Volgens de definitie betekent dit dat $\tuple{a+b, b+a} \sim_\Int
\tuple{0,0}$, en dus dat $\equivrep{a+b, b+a}{} = \equivrep{0, 0}{}
= 0_\Int$. We krijgen daarom $i + j =
\equivrep{a,b}{}+\equivrep{b,a}{}=\equivrep{a+b, b+a}{}= \equivrep{0,
0}{} = 0_\Int$.
\end{proof}
Hierna bewijzen we de \emph{distributiviteit}: als je met een som
vermenigvuldigt, mag je de vermenigvuldiging over beide termen
verdelen.
\begin{proof}
Kies net als hierboven $i = \equivrep{a_1, b_1}{}$ en $j =
\equivrep{a_2,b_2}{}$ en $k = \equivrep{a_3, b_3}{}$. Dan rekenen we:
\begin{align*}
	i \times (j + k) 
	&= \equivrep{a_1, b_1}{} \times (\equivrep{a_2,b_2}{} + \equivrep{a_3, b_3}{})\\
	&= \equivrep{a_1, b_1}{} \times \equivrep{a_2 + a_3,b_2+b_3}{}\\
	&= \equivrep{a_1  (a_2 + a_3) + b_1  (b_2+b_3), a_1  (b_2 + b_3) + b_1 (a_2 + a_3)}{}\\
	&= \equivrep{a_1 a_2 + a_1a_3 + b_1 b_2+b_1b_3, a_1 b_2 + a_1b_3 + a_2b_1 + a_3b_1}{}\\		
%		&= \equivrep{a_1 a_2 + b_1b_2 + a_1a_3 + b_1 b_2+b_1b_3, a_1 b_2 + a_1b_3 + a_2b_1 + a_3b_1}{}\\		
	&= \equivrep{a_1a_2 + b_1b_2, a_1b_2 + a_2b_1}{} + \equivrep{a_1a_3 + b_1b_3, a_1b_3 + a_3b_1}{}\\
	&= (\equivrep{a_1, b_1}{} \times \equivrep{a_2,b_2}{}) + (\equivrep{a_1, b_1}{} \times  \equivrep{a_3, b_3}{})\\
	&= (i \times j) + (i \times k)
\end{align*}
\end{proof}
De andere vijf regels mag je zelf bewijzen. Als dat is gedaan, weten
we dat $\Int$ een commutatieve ring is. Dus optellen en
vermenigvuldigen werken, met onze definities, zoals ze moeten werken.
\begin{prob}
Bewijs dat $\Int$ een commutatieve ring is.
\end{prob}
Maar we zijn nog niet klaar. We hebben niet alleen vastgelegd hoe
optellen en vermenigvuldigen op $\Int$ werken, maar ook wanneer het
ene getal kleiner dan of gelijk aan het andere is; daarvoor gebruiken
we de ordeningsrelatie $\leq$. Ook die moet werken zoals bedoeld.
Preciezer gezegd moeten we bewijzen dat $\Int$ een \emph{geordende}
ring is.\footnote{Bedenk dat volgens
	\olref[sfr][rel][ord]{def:linearorder} een totale ordening
	een relatie is die reflexief, transitief, antisymmetrisch en samenhangend is.
	In een ordening heet die laatste eigenschap soms
	\emph{trichotomie}. Ze zegt dat voor elke $a$ en $b$ geldt: $a \leq b \lor a
	= b \lor a \geq b$.} 
\begin{defn}
Een \emph{geordende ring} is een commutatieve ring met bovendien een
totale ordening $\leq$. Daarbij moeten deze twee regels gelden:
\begin{align*}
	a \leq b &\lif a + c \leq b + c\\
	(a \leq b \land 0 \leq c) &\lif a \times c \leq b \times c
\end{align*}
\end{defn}
\begin{prob}
Bewijs dat $\Int$ een geordende ring is. 
\end{prob}
Ook dit bewijs kost veel schrijfwerk, maar de stappen zelf zijn
eenvoudig. We laten het aan jou over om te bewijzen dat onze
$\Int$ inderdaad een geordende ring is.
Daarmee zijn we klaar met de gehele getallen. Voor de rationale getallen moeten we
bijna hetzelfde doen. We moeten bewijzen dat ze een geordend
\emph{lichaam} vormen, met onze definities van $+$, $\times$ en
$\leq$:
\begin{defn}\ollabel{orderedfield}
Een \emph{geordend lichaam} is een geordende ring waarvoor ook de volgende regel geldt:
\begin{align*}
	\emph{Multiplicatieve inverse (omgekeerde bij vermenigvuldigen)}& & (\forall a \in S \setminus \{0\})(\exists b \in S) a\times b& = 1
\end{align*}
\end{defn}
Als je eenmaal hebt bewezen dat $\Int$ een geordende ring is, is de
volgende stap niet moeilijk. Je moet alleen veel regels nalopen om
te bewijzen dat $\Rat$ een geordend lichaam is.
\begin{prob}
Bewijs dat $\Rat$ een geordend lichaam is.
\end{prob}
Na de gehele en rationale getallen zijn alleen de reële getallen nog
over. We moeten bewijzen dat $\Real$ een \emph{volledig} geordend
lichaam is. Dat is een geordend lichaam met de
volledigheidseigenschap. In \olref[cuts]{realcompleteness} hebben we
al bewezen dat $\Real$ die eigenschap heeft. We moeten nog wel alle
(saaie) stappen nalopen om te bewijzen dat $\Real$ een geordend
lichaam is.
Voor we aan \emph{die} hoop werk beginnen, moeten we eerst een paar
dingen controleren die ``meteen duidelijk'' lijken. We moeten
bijvoorbeeld zeker weten dat $\alpha + \beta$ volgens onze definitie
zelf weer een \emph{snede} is, voor alle sneden $\alpha$ en $\beta$.
Hier is het bewijs:
\begin{proof}	
Omdat $\alpha$ en $\beta$ allebei sneden zijn, is $\alpha + \beta = \Setabs{p
+ q}{p \in \alpha \land q \in \beta}$ een deelverzameling van $\Rat$ die
niet leeg en niet de hele verzameling is. Stel nu dat $x < p + q$ voor een
$p \in \alpha$ en een $q \in \beta$. Dan is $x - p < q$, zodat $x - p
\in \beta$, en dus $x = p + (x - p) \in \alpha + \beta$. Zodra een
getal in deze verzameling zit, zitten alle kleinere rationale
getallen er dus ook in. Dat wil zeggen dat $\alpha + \beta$ een
beginstuk van $\Rat$ is. Neem tot slot een willekeurig $p + q \in
\alpha + \beta$. Omdat $\alpha$ en $\beta$ allebei sneden zijn, kunnen
we $p_1 \in \alpha$ en $q_1 \in \beta$ vinden met $p < p_1$ en $q <
q_1$. Dan is $p + q < p_1 + q_1 \in \alpha + \beta$. Er zit dus
altijd nog een groter getal in $\alpha + \beta$, zodat deze
verzameling geen grootste element, en dus geen maximum, heeft.
\end{proof}
Met vergelijkbare berekeningen kun je laten zien dat $\alpha - \beta$,
$\alpha \times \beta$ en $\alpha \div \beta$ ook sneden zijn. Bij
delingen moet je het geval waarin $\beta$ de nulsnede is buiten
beschouwing laten. Ook dat bewijs laten we aan jou over.
\begin{prob}
Bewijs dat $\Real$ een geordend lichaam is.
\end{prob}
Er blijft nog één klein punt over. In \olref[cuts]{sec} zeggen we dat
we $\sqrt{2} = \Setabs{p \in \Rat}{p < 0 \text{ of }p^2 < 2}$ kunnen nemen.
We moeten nog bewijzen dat die verzameling echt een \emph{snede} is.
Dat gaat zo:
\begin{proof}
Deze verzameling is duidelijk niet leeg, niet gelijk aan alle
rationale getallen en een beginstuk daarvan. We hoeven dus alleen te
bewijzen dat ze geen grootste element heeft. Neem daarvoor een
willekeurig positief rationaal getal $p$ met $p^2 < 2$ en zet $q =
\frac{2p+2}{p+2}$. Het is genoeg om te laten zien dat zowel $p < q$
als $q^2 < 2$ geldt. Voor $p < q$ rekenen we:
\begin{align*}
	p^2 &< 2\\
	p^2 + 2p &< 2 + 2p\\
	p(p + 2) &< 2 + 2p\\
	p &< \tfrac{2+2p}{p+2} = q
\end{align*}
Voor $q^2 < 2$ rekenen we:
\begin{align*}
	p^2 &< 2\\
	2p^2 + 4p + 2 &< p^2 + 4p+ 4\\
	4p^2 + 8p + 4 &< 2(p^2 + 4p + 4)\\
	(2p+2)^2 & <2(p+2)^2\\
	\tfrac{(2p+2)^2}{(p+2)^2} &< 2\\
	q^2 &< 2
\end{align*}
\end{proof}
\end{document}